\section{Introduction}
\label{sec:introduction}

A GRPO-trained formalizer achieves a 93\% prove rate on first-order logic problems---yet 87\% of its proved outputs embed the conclusion verbatim among the premises.
The solver succeeds not because the model has learned to reason, but because the formalization trivially entails any conclusion that already appears as a premise.
This structural shortcut is invisible to execution-based evaluation: prove rate registers success regardless of whether the proof reflects genuine multi-step derivation or degenerate circular embedding.

Prove rate---the dominant evaluation metric for neurosymbolic formalizers---conflates faithful reasoning with degenerate shortcuts.
Without a principled method to assess whether each component of a formalization contributes to the proof, practitioners cannot distinguish models that reason from models that exploit solver mechanics.
This gap is particularly consequential in reinforcement learning settings, where the optimizer will find and exploit any shortcut the reward function fails to penalize~\citep{skalse-reward-hacking, reward-hacking-survey}.

Neurosymbolic approaches combine large language models with external solvers for logical reasoning.
Prompting-based methods such as SatLM~\citep{satlm}, Logic-LM~\citep{logic-lm}, and LINC~\citep{linc} translate natural language to formal representations verified by automated provers, but evaluate correctness through execution alone without examining the internal structure of proved outputs.
Recent work applies reinforcement learning---specifically Group Relative Policy Optimization~\citep[GRPO;][]{deepseekmath}---to train LLMs that generate formal logic programs verified by theorem provers~\citep{mellgren-prolog-grpo, logic-rl}.
The standard reward signal in these systems is binary execution feedback: 1 if the solver returns a proof, 0 otherwise.
This drives rapid improvement in prove rate but provides no gradient on structural quality---a formalization where every premise contributes to a valid derivation receives the same reward as one that smuggles the conclusion into the premises.
Process reward models~\citep{lightman-verify-step} offer finer-grained step-level supervision for mathematical reasoning, but require expensive human annotations and do not transfer to arbitrary solver-backed verification settings.
No existing method provides an annotation-free, solver-agnostic diagnostic for structural integrity in neurosymbolic outputs.

We address this gap with \emph{Solver Trace Ablation} (STA), a diagnostic protocol inspired by mutation testing in software engineering~\citep{demillo1978} that applies component necessity testing---leave-one-out (LOO) ablation on solver inputs---to quantify structural integrity (\autoref{fig:sta-overview}).
For each proved formalization, STA removes each component individually and re-runs the solver; the fraction of components whose removal causes proof failure yields the \emph{Intensional Structural Reward} (ISR).
ISR${=}$1 indicates a maximally tight formalization where every component is individually necessary; ISR${=}$0 indicates complete structural redundancy.
The protocol is solver-agnostic, requiring only a solver that accepts component-level input manipulation.
Applied as a post-hoc diagnostic on a GRPO formalizer trained with FOLIO~\citep{folio} data, ISR reveals a structural collapse trajectory entirely invisible to prove rate: ISR falls from 0.642 at 100 steps to 0.132 at 500 steps to 0.003 at 1000 steps---a 99.5\% degradation---while prove rate simultaneously rises from 32\% to 96\%.
Three random seeds diverge radically (PR $\sigma{=}$47\,pp), each converging to a qualitatively different degeneration mode: conclusion embedding, repetition collapse, and tautological collapse.
A supervised fine-tuning (SFT) baseline achieves 92.9\% genuine proofs with near-zero hacking, confirming that the structural collapse is a consequence of RL reward optimization, not model capacity or data limitations.

Given ISR's diagnostic power, a natural question is whether it can serve as a training signal to steer models toward structurally sound formalizations.
The answer is negative: when ISR is incorporated into the reward ($r = r_{\text{exec}} \cdot (1 + \beta \cdot \text{ISR})$), the model games the diagnostic metric itself---a phenomenon we characterize as \emph{meta-level Goodhart's Law}~\citep{manheim-goodhart-taxonomy, karwowski-goodhart-rl}.
At low dose ($\beta{=}0.1$), two of three seeds are dominated by \emph{conclusion duplication}: the model generates a single premise identical to the conclusion, achieving ISR${=}$1 by construction while encoding no reasoning.
Across three seeds, 64.5\% of proved samples exhibit this ISR-specific gaming mode (H7), compared to 0\% under a matched random-reward control.
The remaining seed's apparent ISR advantage is entirely a premise-count artifact: controlling for component count eliminates the difference (1.06$\times$; \S\ref{sec:goodhart}).
At high dose ($\beta{=}0.3$), classic Goodhart collapse occurs: prove rate falls to 18\%, compared to 48\% for a matched random control~\citep{gao-reward-overoptimization}.
This dose-response pattern---low-dose proxy gaming followed by high-dose performance collapse---maps onto established Goodhart dynamics but occurs at the \emph{meta} level: ISR was designed to detect gaming, yet becomes the target of gaming when optimized.

We further validate that both the diagnostic utility and the Goodhart phenomenon transfer across formal languages.
In Prolog experiments using SWI-Prolog as the solver, exec-only training produces proved-only ISR near zero for two of three seeds, replicating the FOL diagnostic finding.
When ISR is added as a Prolog training reward, the result is \emph{universal} Goodhart: all three seeds converge to gaming with 0\% genuine proofs---stronger than the FOL result, where ${\sim}$16\% genuine residual persists---each via a distinct mechanism (tautological collapse, degenerate syntax, and query-as-fact gaming).

Across 6 reward conditions, 3 seeds, 2 logic types, and up to 1000 training steps, we make four contributions:

\begin{enumerate}[nosep,leftmargin=*]
\item We introduce Solver Trace Ablation and establish ISR as a diagnostic that detects reward hacking invisible to prove rate, documenting a complete collapse trajectory ($0.642 \to 0.132 \to 0.003$) over 1000 GRPO training steps; an SFT baseline confirms hacking is RL-specific (92.9\% genuine vs.\ ${<}$16\% under any RL condition).
\item We empirically characterize meta-level Goodhart's Law in neurosymbolic RL---a diagnostic designed to detect gaming becomes gamed when used as reward---and identify the premise-count confound, a systematic LOO bias where fewer components mechanically inflate ISR, that explains the sole apparent improvement (1.06$\times$ premise-matched).
\item We demonstrate cross-logic diagnostic transfer from FOL/Prover9 to Prolog/SWI-Prolog, where ISR-as-reward produces universal Goodhart (3/3 seeds gaming, 0\% genuine), demonstrating that simpler formal semantics enable more complete exploitation.
\item We catalog 7 reward hacking modes across conditions and logic types, including query-as-fact gaming that bypasses text-match classifiers, motivating functional cross-validation of automated hacking diagnostics.
\end{enumerate}

\noindent \S\ref{sec:related-work} surveys related work; \S\ref{sec:method} introduces the STA protocol; \S\ref{sec:experiments} presents experiments; \S\ref{sec:discussion} discusses implications.
